%═══════════════════════════════════════════════════════════════════════════════
\chapter{The Electric Potential}
%═══════════════════════════════════════════════════════════════════════════════

\textit{(Chapter 2, Sections 2.1--2.8)}

%───────────────────────────────────────────────────────────────────────────────
\section{Line Integral of the Electric Field}
%───────────────────────────────────────────────────────────────────────────────

Let $\vect{E}$ be some electric field and $P_1$, $P_2$ some points within it. Then the \emph{line integral} between these points is
\[
  W_{P_1}^{P_2} = \int_{P_1}^{P_2} \vect{E} \cdot d\vect{s}
\]

\begin{figure}[H]
\begin{center}
\begin{tikzpicture}[scale=1.0]
  % Background field lines (rightward)
  \foreach \y in {-1.2,-0.6,0,0.6,1.2} {
    \draw[-{Stealth[length=4pt]}, gray!50] (-2.5,\y) -- (2.5,\y);
  }
  % Points
  \filldraw[black] (-1.5,0) circle (3pt) node[below left]{$P_1$};
  \filldraw[black] (1.5,0) circle (3pt) node[below right]{$P_2$};
  % Path 1: straight
  \draw[thick, blue!70!black, -{Stealth[length=5pt]}]
    (-1.5,0) -- (1.5,0) node[midway, below=4pt]{Path 1};
  % Path 2: curved above
  \draw[thick, red!70!black, -{Stealth[length=5pt]}]
    (-1.5,0) .. controls (-0.5,1.2) and (0.5,1.2) .. (1.5,0)
    node[midway, above=4pt]{Path 2};
  % Label: same integral
  \node[font=\small] at (0,-1.8)
    {$\displaystyle\int_{\text{Path 1}}\vect{E}\cdot d\vect{s} = \int_{\text{Path 2}}\vect{E}\cdot d\vect{s}$};
\end{tikzpicture}
\end{center}
\caption{Two different paths from $P_1$ to $P_2$ through an electric field. Because $\vect{E}$ is conservative ($\curl\vect{E}=0$), the line integral is path-independent.}
\end{figure}

\begin{itemize}
  \item Since the electric field is \emph{conservative}, the line integral is path-independent for all paths that begin the same and end the same.

  \begin{fact}[Corollary]
    $\displaystyle\oint \vect{E} \cdot d\vect{s} = 0$ around any closed path.
  \end{fact}

  \item The line integral is \emph{linear}: if $\vect{E} = \vect{E}_1 + \vect{E}_2 + \dots + \vect{E}_n$, then
  \[
    \int_{P_1}^{P_2} \vect{E} \cdot d\vect{s}
    = \int_{P_1}^{P_2} \vect{E}_1 \cdot d\vect{s}
    + \int_{P_1}^{P_2} \vect{E}_2 \cdot d\vect{s}
    + \dots
    + \int_{P_1}^{P_2} \vect{E}_n \cdot d\vect{s}
  \]

  \item You will never have an electric field that points in a circle.
\end{itemize}

%───────────────────────────────────────────────────────────────────────────────
\section{Potential Difference and the Potential Function}
%───────────────────────────────────────────────────────────────────────────────

\begin{formula}[Electric potential difference]
\[
  \phi_{21} = -\int_{P_1}^{P_2} \vect{E} \cdot d\vect{s}
\]
Work/unit charge in moving a positive charge from $P_1$ to $P_2$, measured in volts $= \frac{1\,\text{J}}{1\,\text{C}}$ (also: $1\,\text{V} = 1\,\frac{\text{J}}{\text{C}}$, ``how hard it is to push 1 coulomb through'').
\end{formula}

\begin{itemize}
  \item If you fix one point, $P_1$, the $\phi$ becomes a scalar function of $P_2$ $(x,y,z)$: $\phi(x,y,z)$.
\end{itemize}

\begin{formula}[Potential from a point charge at the origin]
\[
  \phi(r) = -\frac{1}{4\pi\varepsilon_0}\int_{\infty}^{r} \frac{Q}{r'^2}\,dr'
           = \frac{1}{4\pi\varepsilon_0}\frac{Q}{r}
\]
\textit{Note: there is no negative sign!}
\end{formula}

\begin{itemize}
  \item Don't confuse $\phi(r)$ with $U$: $\phi(r)$ is the potential function; $U$ is the total work required to assemble a system (not to move a test charge).
\end{itemize}

%───────────────────────────────────────────────────────────────────────────────
\subsection*{Example: Potential Due to a Uniform Sphere}
%───────────────────────────────────────────────────────────────────────────────

\begin{example}[Potential of a uniformly charged sphere (radius $R$, volume charge density $\rho$)]
Goal: find the potential for all values of $r$, both inside and outside the sphere.

\begin{itemize}
  \item Inside the sphere: $E(r) = \dfrac{\rho r}{3\varepsilon_0}$
  \item Outside: $E(r) = \dfrac{\rho R^3}{3\varepsilon_0 r^2}$
\end{itemize}

Thus,
\[
  \phi_\text{out}(r) = -\int_\infty^r E(r')\,dr'
  = -\int_\infty^r \frac{\rho R^3}{3\varepsilon_0 r'^2}\,dr'
  = \frac{\rho R^3}{3\varepsilon_0 r}
\]

In terms of the total charge $Q = \frac{4}{3}\pi R^3 \rho$:
\[
  \phi_\text{out}(r) = \frac{Q}{4\pi\varepsilon_0 r} \quad \text{(as expected)}
\]

For $\phi_\text{in}$:
\begin{align*}
  \phi_\text{in}(r)
  &= -\int_\infty^R E(r')\,dr' - \int_R^r E(r')\,dr' \\
  &= -\int_\infty^R \frac{\rho R^3}{3\varepsilon_0 r'^2}\,dr'
     - \int_R^r \frac{\rho r'}{3\varepsilon_0}\,dr' \\
  &= \frac{\rho R^3}{3\varepsilon_0 R} - \frac{\rho}{6\varepsilon_0}(r^2 - R^2)
   = \frac{\rho R^2}{\varepsilon_0} \cdot \frac{1}{2\varepsilon_0} - \frac{\rho r^2}{6\varepsilon_0}
   = \frac{\rho R^2}{2\varepsilon_0} - \frac{\rho r^2}{6\varepsilon_0}
\end{align*}
\end{example}

%───────────────────────────────────────────────────────────────────────────────
\section{Gradient of a Scalar Function}
%───────────────────────────────────────────────────────────────────────────────

\textit{(Linear small-volume geometry: Van der Graaff)}

Using the potential function we can also go the other way and find the electric field:

\begin{formula}[Gradient in Cartesian coordinates]
\[
  \grad f = \hat{x}\,\pdv{f}{x} + \hat{y}\,\pdv{f}{y} + \hat{z}\,\pdv{f}{z}
\]
\end{formula}

\begin{itemize}
  \item Tells how the function $f$ varies in a small neighbourhood.
  \item The direction of the gradient is in the direction of steepest ascent.
  \item If $f$ is a function of $r$ only, then the direction of steepest ascent is always along the radial direction.
\end{itemize}

%───────────────────────────────────────────────────────────────────────────────
\subsection{Derivation of the Field from the Potential}
%───────────────────────────────────────────────────────────────────────────────

Consider the value of $\phi$ at two nearby points $(x,y,z)$ and $(x+dx,\,y+dy,\,z+dz)$:
\[
  d\phi = \pdv{\phi}{x}\,dx + \pdv{\phi}{y}\,dy + \pdv{\phi}{z}\,dz
  \quad \text{(to first order)}
\]

But $d\phi = -\vect{E}\cdot d\vect{s} = -\vect{E}\cdot(\hat{x}\,dx + \hat{y}\,dy + \hat{z}\,dz)$, thus:
\[
  E_x = -\pdv{\phi}{x}, \quad E_y = -\pdv{\phi}{y}, \quad E_z = -\pdv{\phi}{z}
\]

\begin{formula}[Field from potential]
\[
  \boxed{\vect{E} = -\grad\phi}
\]
\end{formula}

%───────────────────────────────────────────────────────────────────────────────
\subsection{Equipotential Lines}
%───────────────────────────────────────────────────────────────────────────────

The condition for equipotentials is:
\[
  dx\,\pdv{\phi}{x} + dy\,\pdv{\phi}{y} + dz\,\pdv{\phi}{z} = 0
  \implies \left(\pdv{\phi}{x}, \pdv{\phi}{y}\right)\cdot(dx,dy) = 0
\]

\begin{figure}[H]
\begin{center}
\begin{tikzpicture}[scale=1.0]
  % Equipotential circles (dashed, blue)
  \foreach \r in {0.6,1.0,1.4,1.9} {
    \draw[dashed, blue!60!black, thick] (0,0) circle (\r);
    \node[blue!60!black, font=\tiny] at ({(\r+0.15)*cos(20)},{(\r+0.15)*sin(20)}) {$\phi\!=\text{const}$};
  }
  % Radial field lines (red)
  \foreach \angle in {0,45,90,135,180,225,270,315} {
    \draw[-{Stealth[length=4pt]}, red!70!black, thick]
      ({0.45*cos(\angle)},{0.45*sin(\angle)}) -- ({1.9*cos(\angle)},{1.9*sin(\angle)});
  }
  % Charge at center
  \filldraw[black] (0,0) circle (3.5pt) node[below right, font=\small]{$+q$};
  % Labels
  \node[red!70!black, font=\small] at (2.5,-0.3) {$\vect{E}$ (radial)};
  \node[blue!60!black, font=\small] at (-2.5,0.3) {equip.\ circles};
\end{tikzpicture}
\end{center}
\caption{Equipotential surfaces (dashed circles, $\phi = Q/4\pi\varepsilon_0 r = \text{const}$) and radial electric field lines for a positive point charge. Field lines are always perpendicular to equipotentials.}
\end{figure}

\textbf{Methods to find equipotential lines:}
\begin{enumerate}
  \item Find equipotential lines directly from $\phi$.
  \item Find $\dfrac{dy}{dx}$, then invert it to $-\dfrac{\partial\phi/\partial x}{\partial\phi/\partial y}$.
  \item Integrate: $\dfrac{dy}{dx}$. \textbf{OR}
  \item Set $\dfrac{dy}{dx} = \dfrac{E_y}{E_x}$ at every point if you have the field lines.
\end{enumerate}

\textbf{Finding field lines} (not equipotential lines): the ratios $dx:dy:dz = E_x:E_y:E_z$.

\paragraph{Example:} $\phi(x,y) = 6\ln(x) - 6\ln(y)$

\[
  \frac{dy}{dx} = \frac{E_y}{E_x}
  = \frac{6\ln(y)\cos(y)}{6\ln(y)\cos(x)} = \frac{\tan(x)}{\tan(y)}
\]
\[
  \tan(y)\,dy = \tan(x)\,dx
  \implies \ln(\cos(x)) = \ln(\cos(y)) + C
  \implies y = \cos^{-1}(A\cos(x))
\]

%───────────────────────────────────────────────────────────────────────────────
\subsection{Potential of a Charge Distribution}
%───────────────────────────────────────────────────────────────────────────────

\begin{itemize}
  \item The potential at any isolated point is $\dfrac{q}{4\pi\varepsilon_0 r}$.
  \item Superposition works for potentials as well as fields.
  \item If all the sources are confined in some finite region, it is simplest to put the zero of potential at infinite distance.
\end{itemize}

\begin{formula}[Potential from a continuous charge distribution]
\[
  \phi(x,y,z) = \int_\text{source} \frac{\rho(x',y',z')\,dx'\,dy'\,dz'}{4\pi\varepsilon_0\, r}
\]
where $r$ is the distance from the point $(x',y',z')$ (which the potential is being evaluated from) to $(x,y,z)$.
\end{formula}

\[
  \phi_0 = \phi_0 + \phi_1 + \cdots
\]

\begin{itemize}
  \item One caveat: for infinite systems such as an infinitely long line charge, the integral diverges.
\end{itemize}

%───────────────────────────────────────────────────────────────────────────────
\subsection*{Example: Potential of a Long Charged Wire}
%───────────────────────────────────────────────────────────────────────────────

\begin{example}[Potential difference for an infinite line charge (linear charge density $\lambda$)]
The field is $E(r) = \dfrac{\lambda}{2\pi\varepsilon_0 r}$.
\begin{align*}
  \phi_{21}
  &= -\int_{P_1}^{P_2} \vect{E}\cdot d\vect{s}
   = -\int_{r_1}^{r_2} \frac{\lambda}{2\pi\varepsilon_0 r}\,dr \\
  &= -\frac{\lambda}{2\pi\varepsilon_0}\bigl(\ln(r_2) - \ln(r_1)\bigr) \\
  &= -\frac{\lambda}{2\pi\varepsilon_0}\ln(r_2) + \text{constant}
\end{align*}

To verify: $\vect{E}(r) = -\grad\phi = -\dfrac{\lambda}{2\pi\varepsilon_0}\left(\pdv{\ln r}{x}, \pdv{\ln r}{y}, \pdv{\ln r}{z}\right)$

\textbf{Note: blows up.}

\[
  = \frac{\lambda}{2\pi\varepsilon_0}\left(\frac{x}{r^2}, \frac{y}{r^2}, \frac{z}{r^2}\right)
  = \frac{\lambda}{2\pi\varepsilon_0 r}\,\hat{r}
  = \frac{\lambda}{2\pi\varepsilon_0}\hat{r}
\]
\end{example}

%───────────────────────────────────────────────────────────────────────────────
\section{Uniformly Charged Disk}
%───────────────────────────────────────────────────────────────────────────────

\begin{figure}[H]
\centering
\caption{Uniformly charged disk (surface charge density $\sigma$, radius $a$) with a field point $P_1$ on the axis at distance $y$.}
\end{figure}

Area of the ring segment:
\[
  dA = 2\pi s \, ds \implies d\phi_\text{ring} = \frac{\sigma\,dA}{2\pi\varepsilon_0\sqrt{y^2+s^2}}
     = \frac{\sigma\,s\,ds}{2\varepsilon_0\sqrt{y^2+s^2}}
\]

Since each point is a distance $\sqrt{y^2+s^2}$ from each point on the ring, the contribution of the ring to the potential is $d\phi = \dfrac{\sigma\,s\,ds}{2\varepsilon_0\sqrt{y^2+s^2}}$.

Thus,
\[
  \phi(0,y,0) = \int_0^a \frac{\sigma\,s\,ds}{2\varepsilon_0\sqrt{y^2+s^2}}
  = \frac{\sigma}{2\varepsilon_0}\sqrt{y^2+s^2}\Big|_0^a
\]

\begin{formula}[Potential on the axis of a uniformly charged disk]
\[
  \phi(0,y,0) = \frac{\sigma}{2\varepsilon_0}\left(\sqrt{y^2+a^2} - y\right)
  \quad \text{for } y > 0
\]
\end{formula}

When $y=0$, there is a discontinuity in the graph of $\phi(y)$, so $\phi(0) = \dfrac{\sigma a}{2\varepsilon_0}$.

When $y \gg a$, we can approximate $\sqrt{y^2+a^2} - y$ as:
\[
  \sqrt{y^2+a^2} - y
  = y\left[\left(1+\frac{a^2}{y^2}\right)^{1/2} - 1\right]
  = y\left[1 + \tfrac{1}{2}\left(\frac{a^2}{y^2}\right) + \dots - 1\right]
  \approx \frac{a^2}{2y}
\]
Thus,
\[
  \phi(0,y,0) \approx \frac{\pi a^2 \sigma}{4\pi\varepsilon_0 y}
  \quad \text{for } y \gg a
\]
which is the same as the expression for a point charge of the same magnitude.

%───────────────────────────────────────────────────────────────────────────────
\section{Dipoles}
%───────────────────────────────────────────────────────────────────────────────

Consider a setup with 2 equal/opposite charges; the test charge is at $r \gg \ell$.

\begin{figure}[H]
\centering
\caption{Dipole: charge $-q$ at the origin, charge $+q$ at distance $\ell$, field point $P$ at $(r,\theta)$ with $r \gg \ell$.}
\end{figure}

\subsection{Calculation of $\phi$ and \texorpdfstring{$\vect{E}$}{E}}

We will use polar coords ($\theta$, $r$); $\phi$ doesn't matter.

The exact expression for the potential is:
\[
  \phi_p = \frac{kq}{r_+} - \frac{kq}{r_-} \tag{1}
\]

Instead of being exact of course, we will approximate $\hat{r}_+$ and $\hat{r}_-$ as being effectively parallel for $r \gg \ell$ in the limit.

The picture becomes: we can see that $r_+ \approx r - \frac{\ell}{2}\cos\theta$ and $r_- \approx r + \frac{\ell}{2}\cos\theta$.

Using the approximation $\frac{1}{1\pm\delta} \approx 1 \mp \delta$, (1) becomes:

\begin{align*}
  \phi(r,\theta)
  &= \frac{kq}{r - \frac{\ell}{2}\cos\theta} - \frac{kq}{r + \frac{\ell}{2}\cos\theta} \\
  &= \frac{kq}{r}\left(\frac{1}{1-\frac{\ell\cos\theta}{2r}} - \frac{1}{1+\frac{\ell\cos\theta}{2r}}\right) \\
  &= \frac{kq}{r}\left(\left(1 + \frac{\ell\cos\theta}{2r}\right) - \left(1 - \frac{\ell\cos\theta}{2r}\right)\right) \\
  &= \frac{kq\ell\cos\theta}{r^2} = \frac{q\ell\cos\theta}{4\pi\varepsilon_0 r^2}
\end{align*}

\begin{formula}[Dipole potential]
\[
  \phi(r,\theta) = \frac{p\cos\theta}{4\pi\varepsilon_0 r^2}, \quad r \gg \ell
\]
where $p = q\ell$ is the \emph{dipole moment}.
\end{formula}

\begin{itemize}
  \item Note that $\phi(r,\theta) \propto \frac{1}{r^2}$, not $\frac{1}{r}$.
  \item Makes sense as the potentials from the two opposite charges nearly cancel.
  \item There is no angular dependence on $\phi$.
  \item Part of a general technique called \emph{multipole expansion}.
\end{itemize}

\begin{figure}[H]
\centering
\caption{Multipole hierarchy: monopole (single charge), dipole (two opposite charges), bi-monopole, quadrupole.}
\end{figure}

\subsection{Finding \texorpdfstring{$\vect{E}$}{E} in a Dipole --- Derivation of Spherical Gradient}

In spherical coordinates, ($\theta$ and $\phi$ must be kept separate):
\[
  \grad\phi = \pdv{\phi}{r}\,\hat{r} + \frac{1}{r}\pdv{\phi}{\theta}\,\hat{\theta}
\]

Thus,
\begin{align*}
  \vect{E}(r,\theta)
  &= -\grad\phi
   = -\pdv{}{r}\!\left(\frac{p\cos\theta}{4\pi\varepsilon_0 r^2}\right)\hat{r}
     + \frac{1}{r}\pdv{}{\theta}\!\left(\frac{p\cos\theta}{4\pi\varepsilon_0 r^2}\right)\hat{\theta} \\
  &= \frac{p}{4\pi\varepsilon_0 r^3}(2\cos\theta\,\hat{r} + \sin\theta\,\hat{\theta})
\end{align*}

\subsection{Field Lines for \texorpdfstring{$\vect{E}$}{E}}

\begin{figure}[H]
\centering
\caption{Field lines for the dipole electric field (coloured circles represent equipotential lines $\theta = \pi/2$, field closes at small $r$). Equipotential lines are shown at $\theta = \pi/2$ (equatorial plane).}
\end{figure}

\subsection{Deriving Equations for Field Lines and Equipotential Lines}

Let $\phi = \phi_0$ be constant, then:
\[
  \frac{kq\ell\cos\theta}{r^2} = \phi_0
  \implies r^2 = \left(\frac{kq\ell}{\phi_0}\right)\cos\theta
\]
\[
  r = r_0\sqrt{\cos\theta}, \quad r_0 = \sqrt{\frac{kq\ell}{\phi_0}}
\]

To account for $\cos\theta$ being negative in the other half of the dipole, use absolute values:
\begin{formula}[Equipotential curves for a dipole]
\[
  r = r_0\sqrt{|\cos\theta|}, \quad r_0 = \sqrt{\frac{kq\ell}{\phi_0}}
\]
\end{formula}

The slope in polar coordinates is $\dfrac{1}{r}\dfrac{dr}{d\theta}$:
\[
  \frac{1}{r}\frac{dr}{d\theta}
  = \frac{1}{r_0\sqrt{\cos\theta}}\cdot\frac{d}{d\theta}(r_0\sqrt{\cos\theta})
  = \frac{r_0}{r_0\sqrt{\cos\theta}}\cdot\frac{-\sin\theta}{2\sqrt{\cos\theta}}
  = \frac{-\sin\theta}{2\cos\theta}
\]

\begin{itemize}
  \item For $\theta = \pm\pi/2$, slope is $\to\infty$.
\end{itemize}

We can also go from field lines to equipotential lines:
\begin{itemize}
  \item Slope of the tangent $= \dfrac{E_r}{E_\theta} = \dfrac{2\cos\theta}{\sin\theta}$. Why is this the tangent?
  \item To find the equation for the field-line curves, set
    \[
      \frac{1}{r}\frac{dr}{d\theta} = \frac{2\cos\theta}{\sin\theta}
      \implies \int\frac{dr}{r} = \int\frac{2\cos\theta}{\sin\theta}\,d\theta
    \]
    \[
      \ln(r) = 2\ln(\sin\theta) + C
    \]
    \[
      \implies r = r_0 \sin^2\theta
    \]
    where $r_0 = e^C$ is the radius at $\theta = \frac{\pi}{2}$.
\end{itemize}

$\to$ \textit{Do Exercise 2.63 for practice!}

%───────────────────────────────────────────────────────────────────────────────
\section{Divergence of a Vector Function}
%───────────────────────────────────────────────────────────────────────────────

For some arbitrary surface,
\[
  \Phi_\text{net} = \int_S \vect{F}\cdot d\vect{A}
\]
where $\hat{n}$ is normal to the surface and $dA$ is the area of the surface element.

Now imagine dividing the volume into $N$ pieces.
\begin{itemize}
  \item Both volumes share the surface $S_2$: any flux going out of $S_2$ on one side is the flux coming out into the other side, so the internal contributions cancel.
\end{itemize}

Thus,
\[
  \sum_{i=1}^N \int_{S_i} \vect{F}\cdot d\vect{A}_i = \int_S \vect{F}\cdot d\vect{A} = \Phi
\]

\begin{itemize}
  \item As $N\to\infty$ consider the ratio of $\displaystyle\int_{S_i}\vect{F}\cdot d\vect{A}$ and element size $\to 0$.
  \item $\div \vect{F} \defeq \displaystyle\lim_{V_i\to 0} \frac{1}{V_i}\int_{S_i}\vect{F}\cdot d\vect{A}$: the flux out of $V_i$ per unit volume, as $V_i\to 0$.
\end{itemize}

%───────────────────────────────────────────────────────────────────────────────
\section{Gauss's Theorem and the Differential Form of Gauss's Law}
%───────────────────────────────────────────────────────────────────────────────

We can now turn our long back to the full integrals:
\[
  \int_S \vect{F}\cdot d\vect{A}
  = \sum_{i=1}^N \int_{S_i}\vect{F}\cdot d\vect{A}_i
  = \sum_{i=1}^N V_i \left(\frac{\int_{S_i}\vect{F}\cdot d\vect{A}_i}{V_i}\right)
\]

As $N\to\infty$ and $V_i\to 0$, the sum turns into an integral:

\begin{formula}[Gauss's Theorem / Divergence Theorem]
\[
  \int_S \vect{F}\cdot d\vect{A} = \int_V \div\vect{F}\,dV
\]
(when $\vect{F}$ is any vector field over the volume $V$)
\end{formula}

We can also use the fact that:
\[
  \int_S \vect{E}\cdot d\vect{A} = \frac{Q}{\varepsilon_0} = \frac{1}{\varepsilon_0}\int_V \rho\,dV
\]
\[
  \Rightarrow \int_V \div\vect{E}\,dV = \frac{1}{\varepsilon_0}\int_V \rho\,dV
\]

\begin{formula}[Differential (local) form of Gauss's Law]
\[
  \div\vect{E} = \frac{\rho(\vect{r})}{\varepsilon_0}
  \quad\text{the divergence of the force field $= \rho/\varepsilon_0$ at any point}
\]
\end{formula}

\begin{itemize}
  \item Read the example on page 60!
\end{itemize}

%───────────────────────────────────────────────────────────────────────────────
\section{Derivation of Divergence in Cartesian Coordinates}
%───────────────────────────────────────────────────────────────────────────────

Let $\vect{F}$ be a vector function of $x$, $y$, $z$ coordinates. 3 scalar functions: $F_x(x,y,z)$, $F_y(x,y,z)$, $F_z(x,y,z)$.

\begin{figure}[H]
\centering
\caption{Infinitesimal box of dimensions $\Delta x$, $\Delta y$, $\Delta z$ at $(x,y,z)$ used to derive divergence.}
\end{figure}

\begin{itemize}
  \item Net flux through the top and bottom faces (in the $z$-direction):
    \[
      F_{z,\text{top}} - F_{z,\text{bottom}}
    \]
    This can be approximated as $\left(\pdv{F_z}{\Delta z}\right)\Delta z$.
  \item The average value of $F_z$ on the bottom surface of the box (it is constant on the rectangle):
    \[
      \Rightarrow F_{z,\text{bottom}} = F_z(x,y,z) + \frac{\Delta x}{2}\pdv{F_z}{x} + \frac{\Delta y}{2}\pdv{F_z}{y}
    \]
  \item For the top face, the story is similar:
    \[
      F_{z,\text{top}} = F_z(x,y,z+\Delta z) + \frac{\Delta x}{2}\pdv{F_z}{x} + \frac{\Delta y}{2}\pdv{F_z}{y} + \Delta z\pdv{F_z}{\Delta z}
    \]
  \item Thus the net flux out of the box in the $z$-direction is:
    \begin{align*}
      \overrightarrow{\Phi}_z
      &= \Delta x\Delta y\left[F_z(x,y,z+\Delta z) + \frac{\Delta x}{2}\pdv{F_z}{x}
         + \frac{\Delta y}{2}\pdv{F_z}{y} + \Delta z\pdv{F_z}{z}\right] \\
      &\quad - \Delta x\Delta y\left[F_z(x,y,z) + \dots\right] \\
      &= \Delta x\Delta y\,\Delta z\left(\pdv{F_z}{z}\right)
    \end{align*}
\end{itemize}

Thus, for all pairs of sides:
\[
  \Phi = \Delta x\,\Delta y\,\Delta z\left(\pdv{F_x}{x} + \pdv{F_y}{y} + \pdv{F_z}{z}\right)
\]
\[
  \Phi_\text{volume} = \pdv{F_x}{x} + \pdv{F_y}{y} + \pdv{F_z}{z}
\]

This works for arbitrary regions that divide into boxes, like the original argument for Gauss's Divergence Theorem. Thus,

\begin{formula}[Divergence in Cartesian coordinates]
\[
  \div\vect{F} = \pdv{F_x}{x} + \pdv{F_y}{y} + \pdv{F_z}{z} = \nabla\cdot\vect{F} = \frac{\rho(r)}{\varepsilon_0}
\]
\end{formula}

\subsection*{Divergence on a Uniform Sphere}

By Gauss's law outside: $4\pi r^2 E_r = \dfrac{\frac{4}{3}\pi r^3 \rho}{\varepsilon_0}$,
\[
  \Rightarrow \vect{E} = \frac{\rho r}{3\varepsilon_0}\,\hat{r}
\]

For $r > a$, $\hat{r} = \frac{1}{r}(x,y,z)$, so:
\[
  E_x \approx \frac{\rho a^3}{3\varepsilon_0}\cdot\frac{x}{r^3}; \quad E_y, E_z \text{ are similar}
\]

\[
  \pdv{E_x}{x} = \frac{\rho a^3}{3\varepsilon_0}\left(\frac{1}{r^3} - \frac{3x^2}{r^4}\cdot\frac{\partial r}{\partial x}\right)
  = \frac{\rho a^3}{3\varepsilon_0}\left(\frac{1}{r^3} - \frac{3x^2}{r^5}\right)
\]

\[
  \text{thus } {-\div\vect{E}} = \frac{\rho a^3}{3\varepsilon_0}\left(\frac{3}{r^3} - \frac{3}{r^3}\right) = 0
\]

For $r < a$, $E_x = \dfrac{\rho}{3\varepsilon_0} x$, $y$, $z$ --- same:
\[
  \pdv{E_x}{x} = \frac{\rho}{3\varepsilon_0}
  \implies \div\vect{E} = \frac{\rho}{\varepsilon_0}
\]

%───────────────────────────────────────────────────────────────────────────────
\subsection*{Van de Graaff Generator}
%───────────────────────────────────────────────────────────────────────────────

\begin{figure}[H]
\centering
\caption{Van de Graaff generator diagram: charge accumulates on the outer sphere; no charge accumulates inside. The ground reference is at the voltage source.}
\end{figure}

\begin{itemize}
  \item No charge accumulates inside.
  \item If you have a cavity in a metal the electric field inside is zero.
  \item Charge density is higher at a sharp point: used by lightning rods because it discharges charge into the air continuously, ionizing the air.
\end{itemize}

%───────────────────────────────────────────────────────────────────────────────
\section{The Laplacian}
%───────────────────────────────────────────────────────────────────────────────

\begin{align*}
  \div\vect{E} &= -\div(\grad\phi)
  = -\left(\pdv[2]{\phi}{x} + \pdv[2]{\phi}{y} + \pdv[2]{\phi}{z}\right)
\end{align*}

\begin{formula}[Laplacian and Poisson's equation]
\[
  \div\vect{E} = -\laplacian\phi, \quad
  \text{where } \laplacian = \nabla\cdot\nabla
  = \pdv[2]{}{x} + \pdv[2]{}{y} + \pdv[2]{}{z}
\]
Since we know $\div\vect{E} = -\dfrac{\rho}{\varepsilon_0}$:
\[
  \laplacian\phi = -\frac{\rho}{\varepsilon_0}
\]
\end{formula}

\textbf{Useful equations to remember in spherical coordinates:}
\[
  \grad\phi = \frac{1}{r}\pdv{}{r}(r\phi), \qquad
  \laplacian\phi = \frac{1}{r^2}\pdv{}{r}(r\phi)
\]

\subsection*{Example}

Let $\phi(r) = \dfrac{q}{4\pi\varepsilon_0 r}\exp\!\left(-\dfrac{r}{r_0}\right)$.

In spherical coordinates: $\nabla = \left(\pdv{}{r},\, \dfrac{1}{r}\pdv{}{\theta},\, \dfrac{1}{r\sin\theta}\pdv{}{\phi}\right)$.

Thus $\grad\phi = \vect{E} = -\left(\pdv{\phi}{r}, 0, 0\right)$:
\[
  = \frac{q}{4\pi\varepsilon_0 r}\left(\frac{1}{r_0}\exp\!\left(-\frac{r}{r_0}\right)
    + \exp\!\left(-\frac{r}{r_0}\right)\cdot\frac{-1}{r^2}\right)
\]
\[
  = \frac{q}{4\pi\varepsilon_0 r^2}\left[\frac{1}{r^2} - \frac{1}{r_0 r}\right]\exp\!\left(-\frac{r}{r_0}\right)
\]

Working out $\div\vect{E}$ in spherical coordinates using the thin shell $[r, r+dr]$:
\[
  \Delta\Phi_E = 4\pi(r+dr)^2 E_r(r+dr) - 4\pi r^2 E_r(r)
\]
\[
  \Delta\Phi_E = \pdv{}{r}\bigl(4\pi r^2 E_r\bigr)\,dr
\]
Thus,
\[
  \div\vect{E} = \frac{\pdv{}{r}(4\pi r^2 E_r)\,dr}{4\pi r^2\,dr}
  = \frac{1}{r^2}\pdv{}{r}(r^2 E_r)
\]

So:
\[
  r^2 E_r = \frac{q}{4\pi\varepsilon_0}\left(1 + \frac{r}{r_0}\right)\exp\!\left(-\frac{r}{r_0}\right)
\]

\[
  \Rightarrow \rho = \frac{-q}{4\pi r}\left(\frac{1}{r_0^2}\right)\exp\!\left(-\frac{r}{r_0}\right)
\]

\[
  \int_0^\infty \rho(\vec{r})\,4\pi r^2\,dr = -q \quad ??
\]

%───────────────────────────────────────────────────────────────────────────────
\section{Laplace's Equation}
%───────────────────────────────────────────────────────────────────────────────

\begin{formula}[Laplace's equation]
\[
  \laplacian\phi = 0 \quad \text{whenever } \rho = 0, \text{ since there is no charge}
\]
\end{formula}

\begin{itemize}
  \item The class of functions that satisfy Laplace's equation are called \emph{harmonic functions}, and they have the following property:
\end{itemize}

\begin{theorem}[2.1]
If $\phi(x,y,z)$ satisfies Laplace's equation, then the average value of $\phi$ over the surface of any sphere (not necessarily a small sphere) is equal to the value of $\phi$ at the \emph{centre} of the sphere.
\end{theorem}

\begin{proof}
(For the special case for $\phi$ in regions containing no charge.)

Consider a point charge $q$ and a spherical surface $S$ over which a charge $q'$ is uniformly distributed.
\begin{itemize}
  \item Let $q'$ be brought in from infinity to a distance $R$.
  \item Thus, $W = \dfrac{qq'}{4\pi\varepsilon_0 R}$.
\end{itemize}
Now consider if the sphere was brought in from infinity:
\begin{itemize}
  \item Work is the same: $\dfrac{qq'}{4\pi\varepsilon_0 R}$.
  \item Thus, the average value of $q'$ over the surface is $\dfrac{q}{4\pi\varepsilon_0 R}$, which is the same as the potential.
\end{itemize}
This proves the theorem for a single point charge, but by superposition it should also work for multiple charges, since the average of a sum is the sum of the averages.
\end{proof}

\begin{theorem}[Earnshaw's Theorem]
It is impossible to construct an electrostatic field that will hold a charged particle in stable equilibrium in free (empty) space.
\end{theorem}

\begin{proof}
Gauss's Law proof: Suppose we had an electric field in which, contrary to the theorem, there is a point $P$ such that $\phi$ is at a local minimum. But that means $\vect{E}$ must be pointing everywhere inward, so the flux on a small sphere surrounding $P$ is negative $\Rightarrow$ contradicts Gauss's Law (no charge $\Rightarrow$ flux $= 0$!).

\textit{Second method:} A stable position must be one where $\phi$ is at a local minimum $\Rightarrow$ $\phi$ must be lower or equal than $\phi$ at neighbouring points. However, this is impossible when $\phi$ must be the average of its neighbouring ball.
\end{proof}

\subsection*{Distinguishing the Physics from the Mathematics}

\begin{itemize}
  \item Suppose that Coulomb's law is instead $\propto \dfrac{e^{-2r}}{r^2}$, which has an effectively finite range.
  \item Then $\displaystyle\int_S \vect{E}\cdot d\vect{A} \neq \frac{1}{\varepsilon_0}\int_V \rho\,dV$, since the flux would go to 0 as the surface expands --- \textit{see this for yourself!}
  \item However, $\displaystyle\int_S \vect{F}\cdot d\vect{A} = \int_V \div\vect{F}\,dV$ since that is a mathematical law.
  \item $\div\vect{F} \neq \dfrac{\rho}{\varepsilon_0}$ because a region with 0 density of charge could still have a net flux due to the force field having a finite range (at order $\varepsilon_0$) --- some.
  \item Remember that statements like $\laplacian\phi = -\dfrac{\rho}{\varepsilon_0}$ only make sense at a macroscopic scale since electrons are discrete.
\end{itemize}

\subsection*{Equation Graph}

\begin{figure}[H]
\centering
\caption{Relationship diagram between charge density $\rho$, electric field $\vect{E}$, and electric potential $\phi$: $\vect{E} = \frac{1}{4\pi\varepsilon_0}\int\frac{\rho\,\hat{r}}{r^2}\,dV$ (via matter), $\vect{E} = -\grad\phi$, $\nabla\cdot\vect{E} = \div\vect{E} = \frac{\rho}{\varepsilon_0}$ (Coulomb/Gauss), $\phi = -\int\vect{E}\cdot d\vect{s}$, $-\laplacian\phi = \frac{\rho}{\varepsilon_0}$ (Poisson), $\phi = \frac{1}{4\pi\varepsilon_0}\int\frac{\rho}{r}\,dV$.}
\end{figure}

%───────────────────────────────────────────────────────────────────────────────
\section{Vector Curl: Stokes's Theorem}
%───────────────────────────────────────────────────────────────────────────────

Consider a closed path $C$ along some surface $S$.

\begin{figure}[H]
\centering
\caption{Closed path $C$ on surface $S$, divided into two sub-paths by a bridge $B$; shared contributions of the line integrals cancel out.}
\end{figure}

Let $\Gamma = \displaystyle\int_C \vect{F}\cdot d\vect{s}$.

\begin{itemize}
  \item Now bridge $C$ using a new path $B$.
  \item It is easy to see that $\Gamma = \Gamma_1 + \Gamma_2$: the shared contributions of the two line integrals cancel out.
  \item In the limit we infinitely subdivide and consider the ratio of loop circulation to area; but here area is a vector $\hat{a}_i$, normal to $S$.
  \item The limit becomes $\displaystyle\lim_{a_i\to 0}\frac{\Gamma_i}{a_i}$ or $\displaystyle\lim_{a_i\to 0}\frac{\int_{C_i}\vect{F}\cdot d\vect{s}}{a_i}$.
\end{itemize}

\begin{formula}[Curl]
\[
  \curl\vect{F}\cdot\hat{n} = \lim_{a_i\to 0}\frac{\int_{C_i}\vect{F}\cdot d\vect{s}}{a_i}
\]
where $\hat{n}$ is the unit vector normal to the curve.
\end{formula}

\subsection{Stokes's Theorem}

\[
  \Gamma = \int_C \vect{F}\cdot d\vect{s}
  = \sum_{i=1}^N \Gamma_i
  = \sum_{i=1}^N a_i\left(\frac{\Gamma_i}{a_i}\right)
\]

As $N\to\infty$, $\displaystyle\sum_{i=1}^N a_i\!\left(\frac{\Gamma_i}{a_i}\right) = \sum_{i=1}^N a_i(\curl\vect{F})\cdot\hat{n}_i \to \int_S \curl\vect{F}\cdot d\vect{A}$.

Thus we find that:

\begin{formula}[Stokes's Theorem]
\[
  \int_C \vect{F}\cdot d\vect{s} = \int_S \curl\vect{F}\cdot d\vect{A}
\]
\end{formula}

%───────────────────────────────────────────────────────────────────────────────
\subsection{Curl in Cartesian Coordinates}
%───────────────────────────────────────────────────────────────────────────────

\begin{figure}[H]
\centering
\caption{Small rectangle in the $xy$-plane with corners $A$, $B$, $C$, $D$ and dimensions $\Delta x$, $\Delta y$, used to derive the $z$-component of curl.}
\end{figure}

Computing the line integral around the small rectangle:
\begin{align*}
  \int_A^B \vect{F}\cdot d\vect{s}
  &= F_x\!\left(x, y - \frac{\Delta y}{2}\right) \\
  &= \left(F_x(x,y) - \pdv{F_x}{y}\frac{\Delta y}{2}\right)\Delta x
\end{align*}
\begin{align*}
  \int_B^C \vect{F}\cdot d\vect{s}
  &= \left(F_y(x,y) + \pdv{F_y}{x}\frac{\Delta x}{2}\right)\Delta y
\end{align*}
\begin{align*}
  \int_C^D \vect{F}\cdot d\vect{s}
  &= \left(F_z(x,y) + \pdv{F_z}{y}\frac{\Delta y}{2}\right)(-\Delta x)
\end{align*}
\begin{align*}
  \int_D^A \vect{F}\cdot d\vect{s}
  &= \left(F_y(x,y) + \pdv{F_y}{x}\frac{\Delta x}{2}\right)(-\Delta y)
\end{align*}

If you add these together you get:
\[
  \oint_\square \vect{F}\cdot d\vect{s}
  = \Delta x\,\Delta y\left(\pdv{F_y}{x} - \pdv{F_x}{y}\right)
\]

Thus $(\curl\vect{F})_z = \left(\pdv{F_y}{x} - \pdv{F_x}{y}\right) = \nabla\times\vect{F}$

\begin{formula}[Curl in Cartesian coordinates]
\[
  \curl\vect{F}
  = \hat{x}\left(\pdv{F_z}{y} - \pdv{F_y}{z}\right)
  + \hat{y}\left(\pdv{F_x}{z} - \pdv{F_z}{x}\right)
  + \hat{z}\left(\pdv{F_y}{x} - \pdv{F_x}{y}\right)
\]
OR:
\[
  \curl\vect{F} = \nabla\times\vect{F}
  = \begin{vmatrix}\hat{x} & \hat{y} & \hat{z} \\ \partial/\partial x & \partial/\partial y & \partial/\partial z \\ F_x & F_y & F_z\end{vmatrix}
\]
\end{formula}

\paragraph{Example 1:} $\vect{F} = (ky, -kx, 0)$.
\[
  \curl\vect{F} = -2k\,\hat{z}
\]
(The field rotates like a vortex with field lines going: up-right, down-left pattern.)

\paragraph{Example 2:} $\vect{E} = \dfrac{Q}{4\pi\varepsilon_0 r^2}\,\hat{r}$ ($|\vect{r}| < a$):
\[
  \curl\vect{E} =
  \begin{vmatrix}\hat{x} & \hat{y} & \hat{z} \\ \partial/\partial x & \partial/\partial y & \partial/\partial z \\ x & y & z\end{vmatrix}
  = \vec{0}
\]

%───────────────────────────────────────────────────────────────────────────────
\subsection{Physical Meaning of Curl}
%───────────────────────────────────────────────────────────────────────────────

\begin{itemize}
  \item Imagine a velocity vector field $\vect{G}$, and suppose that $\curl\vect{G} \neq 0$.
  \item Behaviours then go like: $\nearrow\swarrow$ or $\uparrow\downarrow$ (shear).
  \item Things would have the tendency to rotate.
  \item But for an \underline{electrostatic} field, curl must be zero because the path integral of any closed loop is 0!
    \begin{itemize}
      \item Thus by Stokes's theorem the total circulation is 0 over any surface $\Rightarrow$ the curl is 0 at every point.
    \end{itemize}
  \item If $\curl\vect{F} = 0$, then $\vect{F}$ is conservative ($\vect{F} = \grad\psi$) or ($\vect{E} = -\grad\phi$).
\end{itemize}

%───────────────────────────────────────────────────────────────────────────────
\subsection{Curl in Spherical Coordinates}
%───────────────────────────────────────────────────────────────────────────────

\begin{formula}[$r$-component of curl in spherical coordinates]
\[
  (\nabla\times\vect{A})_r
  = \frac{1}{r\sin\theta}\left(\pdv{A_\phi\sin\theta}{\theta} - \pdv{A_\theta}{\phi}\right)\hat{r}
\]
\end{formula}

%───────────────────────────────────────────────────────────────────────────────
\section*{Chapter 2 Summary}
%───────────────────────────────────────────────────────────────────────────────

\begin{itemize}
  \item \textbf{Electric potential difference:}
    \[
      \phi_{21} = -\int_{P_1}^{P_2} \vect{E}\cdot d\vect{s}
    \]
    Relative to infinity, the potential due to a charge distribution is:
    \[
      \phi(x,y,z) = \int\frac{\rho(x',y',z')\,dx'\,dy'\,dz'}{4\pi\varepsilon_0 r}
      \quad \text{or} \quad \sum_i \frac{q_i}{4\pi\varepsilon_0 r}
    \]

  \item \textbf{Cartesian gradient:}
    \[
      \grad f = \hat{x}\,\pdv{f}{x} + \hat{y}\,\pdv{f}{y} + \hat{z}\,\pdv{f}{z}
      \quad\text{gives the direction of largest increase}
    \]
    \[
      \vect{E} = -\grad\phi
    \]
    Implies that the lines of the electric field are perpendicular to the surfaces of constant potential.

  \item Electrostatic potential energy difference of a charge $q$ between $P_1$, $P_2$ $= q\phi_{21}$.

  \item Energy required to assemble a group of charges from infinity is:
    \[
      U = \frac{1}{2}\int \rho\phi\,dV
      \quad\text{or}\quad
      \frac{1}{2}\sum_{j=1}^N q_j \sum_{k\neq j}\frac{1}{4\pi\varepsilon_0}\frac{q_k}{r_{jk}}
      \quad\text{(factor of $\frac{1}{2}$ for double counting)}
    \]

  \item \textbf{Dipole:} $\ominus\!\!-\!\!a\!\!-\!\!\oplus$, two charges at distance $a$ apart.
    \begin{itemize}
      \item Dipole moment: $p = q\ell$ (or $p = qa$).
      \item $\phi(r,\theta) = \dfrac{p\cos\theta}{4\pi\varepsilon_0 r^2}$
      \item $\vect{E}(r,\theta) = \dfrac{p}{4\pi\varepsilon_0 r^3}(2\cos\theta\,\hat{r} + \sin\theta\,\hat{\theta})$
    \end{itemize}

  \item \textbf{Divergence:}
    \[
      \div\vect{F} = \nabla\cdot\vect{F} = \pdv{F_x}{x} + \pdv{F_y}{y} + \pdv{F_z}{z}
    \]
    Gauss's Theorem / Divergence Theorem:
    \[
      \int_S \vect{F}\cdot d\vect{A} = \int_V \div\vect{F}\,dV
    \]
    Combining with Gauss's Law gives $\div\vect{E} = \dfrac{\rho}{\varepsilon_0}$.

  \item \textbf{Laplacian} (Cartesian):
    \[
      \laplacian f = \pdv[2]{f}{x} + \pdv[2]{f}{y} + \pdv[2]{f}{z}
    \]
    \[
      \boxed{\laplacian\phi = -\frac{\rho}{\varepsilon_0}} \quad \text{Poisson's equation}
    \]
    Special case: $\laplacian\phi = 0$ $\Rightarrow$ $\phi$ satisfies this equation if it is in a region with no charge density; see Earnshaw's theorem.

  \item \textbf{Curl:}
    \[
      \curl\vect{F}
      = \begin{vmatrix}\hat{x} & \hat{y} & \hat{z} \\ \partial/\partial x & \partial/\partial y & \partial/\partial z \\ F_x & F_y & F_z\end{vmatrix}
    \]
\end{itemize}

%───────────────────────────────────────────────────────────────────────────────
\section*{Recitation 3}
%───────────────────────────────────────────────────────────────────────────────

\begin{figure}[H]
\centering
\caption{Sphere with a spherical hole (of radius $b$) inside. The two spheres share a centre. The field outside the hole region is due to the full sphere minus the hole.}
\end{figure}

\[
  \vect{E} = \frac{\rho}{3\varepsilon_0}\,\vec{r} - \frac{\rho}{3\varepsilon_0}\,\vec{r'} = \frac{\rho}{3\varepsilon_0}(\vec{r}-\vec{r'})
\]
where $\vec{r} = \vec{b} + \vec{r'}$.

\[
  = \frac{\rho}{3\varepsilon_0}\vec{r} - \frac{\rho}{3\varepsilon_0}\vec{r'}
  = \frac{\rho r}{3\varepsilon_0}\hat{r} - \frac{\rho}{3\varepsilon_0}\vec{r'}
  = \frac{\rho}{3\varepsilon_0}\,(\text{constant})
\]

Thus the field inside a sphere with a hole is uniform! It doesn't even depend on the field outside.

\subsection*{Energy of a uniformly charged sphere}

Using $U = \dfrac{1}{2}\varepsilon_0 E^2$:

\begin{align*}
  \mathcal{E}
  &= \frac{1}{2}\varepsilon_0\left(\frac{\rho}{3\varepsilon_0}\right)^2\int_0^a r^2\cdot 4\pi r^2\,dr
   + \frac{1}{2}\varepsilon_0\frac{Q^2}{(4\pi\varepsilon_0)^2}\int_a^\infty \frac{1}{r^4}\cdot 4\pi r^2\,dr \\
  &= \frac{2\pi\rho^2}{9\varepsilon_0}\cdot\frac{a^5}{5} + \frac{Q^2}{8\pi\varepsilon_0}\cdot\frac{1}{a} \\
  &= \frac{2\pi}{9\varepsilon_0}\cdot\frac{Q^2}{16\pi^2 a^6 5}\cdot a^5 + \frac{Q^2}{8\pi\varepsilon_0 a} \\
  &= \frac{Q^2}{40\pi\varepsilon_0 a} + \frac{Q^2}{8\pi\varepsilon_0 a}
   = \frac{6Q^2}{40\pi\varepsilon_0 a}
   = \frac{3}{5}\frac{Q^2}{4\pi\varepsilon_0 a}
\end{align*}

\paragraph{Method 2:}

\begin{figure}[H]
\centering
\caption{Building the uniformly charged sphere shell by shell, each shell of charge $dq$ being brought in from infinity.}
\end{figure}

\[
  dW = \int_{r=0}^{a} \frac{Q\!\left(\frac{r^3}{a^3}\right)}{4\pi\varepsilon_0 r}\cdot\frac{Q}{\frac{4}{3}\pi a^3}\,4\pi r^2\,dr
\]
\[
  = \frac{Q^2}{\frac{4}{3}\pi a^6}\int_{r=0}^{r=a} r^4\,dr
\]
\[
  W = \frac{Q^2}{\frac{4}{3}\pi\varepsilon_0 a^6}\cdot\frac{a^5}{5}
    = \frac{3}{5}\frac{Q^2}{4\pi\varepsilon_0 a}
\]

%───────────────────────────────────────────────────────────────────────────────
\subsection*{Proof of the Divergence Theorem}
%───────────────────────────────────────────────────────────────────────────────

\[
  \div\vect{E} = \frac{\text{flux coming out of a small box } dV}{\text{volume}}
\]

Thus:
\[
  \int_{V_0} \div\vect{E}\,dV = \oint \vect{E}\cdot d\vect{A}
  = \int \frac{\rho(\vec{r})}{\varepsilon_0}\,dV
\]
Thus, $\div\vect{E} = \dfrac{\rho(r)}{\varepsilon_0}$, and $\nabla\cdot\vect{E} = \dfrac{\rho(\vec{r})}{\varepsilon_0}$.

\begin{figure}[H]
\centering
\caption{Flux going into sub-volume 2 equals flux coming out of sub-volume 1 through their shared surface; the internal fluxes cancel in the divergence theorem.}
\end{figure}

In spherical coordinates, for a thin shell $[r, r+dr]$:
\begin{itemize}
  \item Net flux on the ring $= 4\pi(r+dr)^2 E_r(r+dr) - 4\pi r^2 E_r(r)$.
  \item Taking the derivative:
    \[
      \frac{1}{r^2}\pdv{}{r}\!\left(r^2 E_r\right) = \frac{\rho(r)}{\varepsilon_0}
    \]
\end{itemize}
